\chapter{Data and Representation-Learning Methodology}
\label{chap:methods}

% Backward-compatible label in case other files still refer to it.
\label{chap:latent_representations}

This chapter describes the datasets, preprocessing procedures,
representation-learning models, and experimental controls used in
the thesis.

The methodological design follows directly from the research
questions introduced in
Section~\ref{sec:research_questions}.

Particular care is taken to separate model training, feature
selection, and final evaluation so that the number of samples and
markers can be investigated without introducing information leakage.

The feature-attribution algorithms themselves are described in
Chapter~\ref{chap:feature_importance}, while their quantitative
evaluation is deferred to Chapter~\ref{chap:experiments}.


\section{Experimental Design}
\label{sec:experimental_design}

Let

\begin{equation}
    \mathbf{X}
    \in
    \mathbb{R}^{N \times P}
\end{equation}

denote a genotype matrix containing $N$ individuals and $P$ SNP
markers.

Each individual $i$ is associated with a breed label

\begin{equation}
    y_i
    \in
    \{1,\ldots,C\},
\end{equation}

where $C$ is the number of breeds represented in the corresponding
dataset.

The experimental framework varies three principal components:

\begin{enumerate}

    \item the number of individuals available for training within
    each breed;

    \item the subset and number of SNP markers available to the
    downstream classifier;

    \item the method used to rank candidate markers.

\end{enumerate}

The first component is used to investigate sample efficiency
(\textbf{RQ2}), while the second and third address marker efficiency
and marker-selection strategy
(\textbf{RQ3} and \textbf{RQ4}).

Breed-classification performance using the complete genomic
representation provides the reference configuration for
\textbf{RQ1}.

The overall separation between preprocessing, model training, marker
selection, and held-out evaluation is illustrated in
Figure~\ref{fig:experimental_design}.

\begin{figure}[htbp]
\centering
\begin{tikzpicture}[
    >=Latex,
    font=\small,
    box/.style={draw, semithick, rounded corners=2pt, align=center,
                minimum height=10mm, minimum width=39mm, inner sep=4pt, fill=black!2},
    key/.style={box, fill=black!7},
    small/.style={draw, semithick, rounded corners=2pt, align=center,
                  minimum height=9mm, minimum width=34mm, inner sep=3pt, fill=black!2},
    note/.style={draw, dashed, semithick, rounded corners=2pt, align=center,
                 minimum height=9mm, minimum width=27mm, inner sep=3pt},
    arr/.style={->, semithick},
    darr/.style={->, semithick, dashed}
]
\node[key] (raw) at (0,5.0) {Raw genotype data};
\node[box] (qc) at (0,3.65) {Quality control + encoding};
\node[key] (split) at (0,2.3) {Train / validation / test split};
\node[box] (train) at (-2.9,0.85) {Training partition};
\node[box] (test) at (3.2,0.85) {Held-out test partition};

\node[note] (sub) at (-6.15,0.85) {Subsample $n$\\per breed};
\node[small] (repr) at (-3.8,-0.85) {Representation learning};
\node[small] (classic) at (0.0,-0.85) {Classical ranking};
\node[small] (xai) at (-3.8,-2.25) {XAI attribution};
\node[key] (rank) at (-1.9,-3.65) {Training-only SNP ranking};
\node[key] (panel) at (-1.9,-5.05) {Reduced marker panels};
\node[box] (clf) at (-1.9,-6.45) {Downstream classifier};
\node[key] (eval) at (3.2,-6.45) {Final test evaluation};

\draw[arr] (raw)--(qc); \draw[arr] (qc)--(split);
\draw[arr] (split)--(train); \draw[arr] (split)--(test);
\draw[arr] (train)--(repr); \draw[arr] (train)--(classic);
\draw[darr] (train.west) -- (sub.east);
\draw[arr] (repr)--(xai); \draw[arr] (xai)--(rank); \draw[arr] (classic)--(rank);
\draw[arr] (rank)--(panel); \draw[arr] (panel)--(clf); \draw[arr] (clf)--(eval);
\draw[arr] (test) -- (eval);
\end{tikzpicture}

\caption{Leakage-safe experimental design. Model fitting and SNP ranking are performed using training data, whereas the held-out test partition is reserved for final evaluation. The dashed branch represents the repeated per-breed subsampling used in sample-size experiments. Author's elaboration.}
\label{fig:experimental_design}
\end{figure}



% TODO (optional figure): add a leakage-safe experimental workflow
% only if the final diagram is sufficiently readable at thesis-page scale.



\section{Genomic Datasets}
\label{sec:genomic_datasets_preprocessing}

Two livestock genomic datasets are considered.

The primary experiments are conducted on ovine populations, while
an independent caprine dataset is used to investigate whether the
behavior observed in the primary analysis generalizes to another
livestock species.


\subsection{Primary Ovine Dataset}
\label{subsec:ovine_dataset}

% =========================================================================
% TODO: PRIMARY OVINE DATASET SOURCE
%
% The exact provenance of the primary ovine dataset was not
% identifiable from the files supplied so far.
%
% If the dataset is derived from the International Sheep Genomics
% Consortium / Sheep HapMap dataset described by Kijas et al. (2012),
% add:
%
% \cite{kijas2012sheep}
%
% Otherwise replace this comment with the publication/DOI associated
% with the actual ovine dataset used by CREA-ZA.
% =========================================================================

The primary dataset contains SNP genotype information from
\textbf{TODO} sheep breeds and
\textbf{TODO} individuals.

Before quality control, individuals are characterized by
approximately \textbf{TODO} genomic markers.

% TODO: add the ovine breed-distribution table once the exact
% dataset composition has been confirmed.

The ovine dataset is used for the primary breed-classification,
sample-efficiency, and marker-reduction analyses.


\subsection{Caprine Validation Dataset}
\label{subsec:adaptmap_dataset}

The cross-species validation uses genomic data from the ADAPTmap
consortium
\cite{stella2018adaptmap, colli2018genome}.

ADAPTmap was developed to characterize worldwide goat genetic
diversity and population structure and contains genome-wide SNP data
from geographically distributed caprine populations
\cite{stella2018adaptmap, colli2018genome}.

The caprine dataset is used as an independent methodological
validation domain.

The objective is not to test whether identical genomic loci are
discriminative in sheep and goats.
Instead, the validation examines whether similar relationships
between training sample size, marker-panel size, and
classification performance emerge when the same experimental
procedure is applied to a second livestock species.

% TODO: once the Results chapter is finalized, add the exact
% cross-reference to the cross-species validation subsection.


\section{Genotype Quality Control and Encoding}
\label{sec:genotype_preprocessing}

Raw genotype data are subjected to quality-control procedures before
model training.

Common quality-control operations for SNP datasets include
sample-level missingness filtering, marker-level missingness
filtering, Minor Allele Frequency (MAF) filtering, and, when
appropriate for the downstream analysis, Linkage Disequilibrium
(LD) pruning
\cite{purcell2007plink, anderson2010data, chang2015second}.
Previous discriminant analyses in livestock genomics likewise
applied explicit marker- and sample-level quality-control procedures
before multivariate analysis
\cite{dimauro2013canonical, sorbolini2016signatures}.
The thresholds in the present thesis are nevertheless defined by the
actual experimental pipeline and are not inherited automatically
from those earlier studies.

The exact thresholds used in the final experimental pipeline must be
reported explicitly and applied without reference to the final test
labels.

% TODO:
% Replace the following values ONLY with the thresholds actually used:
%
% Sample missingness       = TODO
% Marker missingness       = TODO
% MAF threshold            = TODO
% LD pruning               = TODO / not used
% LD window                = TODO
% LD r^2 threshold         = TODO

A biallelic genotype can be represented using allele dosage

\begin{equation}
    g_{i,j}
    \in
    \{0,1,2\},
    \label{eq:allele_dosage}
\end{equation}

where $g_{i,j}$ denotes the number of alternative alleles observed
for individual $i$ at marker $j$.

For neural architectures in which the genotype state is modeled
categorically, the dosage representation can alternatively be
mapped to a three-channel one-hot vector

\begin{equation}
    g_{i,j}
    \longrightarrow
    \mathbf{x}_{i,j}
    \in
    \{0,1\}^{3}.
    \label{eq:one_hot_genotype}
\end{equation}

The complete encoded genotype tensor is therefore represented as

\begin{equation}
    \mathbf{X}_{\mathrm{tensor}}
    \in
    \mathbb{R}^{N \times P \times 3}.
    \label{eq:genotype_tensor}
\end{equation}

The motivation for using one-hot encoding is not that the
conventional $0/1/2$ dosage representation is mathematically
invalid.
Allele dosage is a standard representation used in statistical
genetics and genome-wide analysis
\cite{purcell2007plink}
and has also been used directly in multivariate livestock genomic
analyses based on discriminant methods
\cite{dimauro2013canonical, sorbolini2016signatures}.

Rather, one-hot encoding is adopted here as a modeling choice that
provides separate channels for the three genotype states and,
consequently, an explicit channel structure for neural
feature-attribution analysis.

% TODO: if both dosage and one-hot encodings are evaluated, add an
% encoding-ablation subsection in the Results chapter and reference it here.
% Otherwise keep only the encoding actually used in the final experiments.


\section{Data Partitioning and Leakage Prevention}
\label{sec:data_partitioning}

Individuals are divided into mutually exclusive training,
validation, and test partitions.

% TODO:
% Insert the actual split used in the final experiments.

The split is stratified by breed whenever allowed by the available
number of individuals.

All model fitting and hyperparameter selection are performed using
the training and validation partitions.

Crucially, marker ranking is also performed without access to the
final test partition.

This includes rankings derived from population statistics,
Random Forest importance, SHAP, Integrated Gradients, or any other
feature-selection procedure subsequently evaluated on the test set.

Feature selection performed using information from the complete
dataset before final evaluation can introduce selection bias and
produce optimistic estimates of generalization performance
\cite{ambroise2002selection}.

The held-out test set is therefore accessed only after model
training, feature ranking, and panel construction have been
completed.


\section{VMGP-Based Variational Representation Model}
\label{sec:vmgp_architecture}

The first representation-learning approach follows the variational
autoencoder paradigm \cite{kingma2013auto,rezende2014stochastic} and
is adapted from the Variational Multi-view Genomics Predictor (VMGP)
architecture \cite{zhao2025vmgp}.

The choice of a variational architecture is motivated by the
high-dimensional structure of genomic data considered in this work.
Each individual is represented by a large number of SNP markers,
while the number of available individuals can be substantially
smaller than the number of genomic variables. Moreover, correlations
among markers may introduce substantial redundancy in the original
input space. A variational bottleneck provides a natural mechanism
for mapping the complete genotype profile into a lower-dimensional,
regularized latent representation.

VMGP was selected as the architectural reference because it applies
variational representation learning specifically to genomic data and
combines latent representation learning with prediction objectives.
The original model was proposed for genomic prediction in plants
\cite{zhao2025vmgp}; therefore, its use in this thesis is not intended
as a direct transfer of an already validated livestock
breed-classification model. Instead, the architecture is adapted to
investigate whether a genomics-oriented variational representation can
retain information relevant to breed discrimination.

In the proposed adaptation, genotype reconstruction and breed
classification provide complementary learning signals. Reconstruction
encourages the latent representation to preserve information about
the complete genomic profile, whereas the classification objective
encourages it to retain information useful for distinguishing breeds.

For an input genotype representation

\begin{equation}
    \mathbf{x}
    \in
    \mathbb{R}^{D},
\end{equation}

an encoder $q_{\phi}$ maps the input to the parameters of an
approximate posterior distribution

\begin{equation}
    q_{\phi}
    (
        \mathbf{z}
        \mid
        \mathbf{x}
    )
    =
    \mathcal{N}
    \left(
        \mathbf{z};
        \boldsymbol{\mu}(\mathbf{x}),
        \operatorname{diag}
        \boldsymbol{\sigma}^{2}(\mathbf{x})
    \right).
    \label{eq:vae_posterior}
\end{equation}

A latent sample is obtained through the reparameterization trick
\cite{kingma2013auto, rezende2014stochastic}

\begin{equation}
    \mathbf{z}
    =
    \boldsymbol{\mu}(\mathbf{x})
    +
    \boldsymbol{\sigma}(\mathbf{x})
    \odot
    \boldsymbol{\epsilon},
    \qquad
    \boldsymbol{\epsilon}
    \sim
    \mathcal{N}
    (
        \mathbf{0},
        \mathbf{I}
    ),
    \label{eq:reparameterization}
\end{equation}

where $\odot$ denotes element-wise multiplication.

The latent representation is decoded to reconstruct the genotype
input.
The VAE objective combines reconstruction error and
Kullback--Leibler regularization
\cite{kingma2013auto}:

\begin{equation}
    \mathcal{L}_{\mathrm{VAE}}
    =
    \mathcal{L}_{\mathrm{rec}}
    +
    \beta
    D_{\mathrm{KL}}
    \left(
        q_{\phi}
        (
            \mathbf{z}
            \mid
            \mathbf{x}
        )
        \parallel
        p(\mathbf{z})
    \right).
    \label{eq:vae_loss}
\end{equation}

For breed-aware representation learning, the architecture additionally
includes a classification objective:

\begin{equation}
    \mathcal{L}_{\mathrm{total}}
    =
    \lambda_{\mathrm{rec}}
    \mathcal{L}_{\mathrm{rec}}
    +
    \lambda_{\mathrm{KL}}
    \mathcal{L}_{\mathrm{KL}}
    +
    \lambda_{\mathrm{cls}}
    \mathcal{L}_{\mathrm{cls}}.
    \label{eq:vmgp_total_loss}
\end{equation}

% TODO:
% Replace the generic formulation above with the EXACT objective
% implemented in your code.
%
% Report:
% - input dimension;
% - hidden layers;
% - latent dimension;
% - activation functions;
% - dropout;
% - batch normalization;
% - decoder structure;
% - classifier head;
% - lambda weights;
% - optimizer;
% - learning rate.
%
% Do NOT place final accuracy here.

% TODO (optional figure): add the final variational architecture diagram
% only after the implementation details below have been frozen.




The implemented variational architecture is summarized in
Figure~\ref{fig:vmgp_architecture}.

\begin{figure}[htbp]
\centering
\begin{tikzpicture}[
    >=Latex,
    font=\small,
    box/.style={draw, semithick, rounded corners=2pt, align=center,
                minimum height=10mm, minimum width=42mm, inner sep=4pt, fill=black!2},
    key/.style={box, fill=black!7},
    small/.style={draw, semithick, rounded corners=2pt, align=center,
                  minimum height=9mm, minimum width=29mm, inner sep=3pt, fill=black!2},
    arr/.style={->, semithick}
]
\node[key] (x) at (0,4.6) {Encoded SNP profile $\mathbf{x}$};
\node[box] (enc) at (0,3.2) {Encoder $q_\phi$};
\node[small] (mu) at (-2.2,1.75) {$\boldsymbol{\mu}(\mathbf{x})$};
\node[small] (lv) at (2.2,1.75) {$\log\boldsymbol{\sigma}^{2}(\mathbf{x})$};
\node[key] (z) at (0,0.15) {Reparameterization $\rightarrow$ latent $\mathbf{z}$};
\node[box] (dec) at (-3.0,-1.35) {Decoder};
\node[box] (head) at (3.0,-1.35) {Breed prediction head};
\node[small] (rec) at (-3.0,-2.8) {Reconstruction};
\node[small] (prob) at (3.0,-2.8) {Breed probabilities};
\node[key] (loss) at (0,-4.35) {Joint objective $\alpha\mathcal{L}_{\mathrm{rec}}+(1-\alpha)\mathcal{L}_{\mathrm{breed}}$, with $\mathcal{L}_{\mathrm{rec}}=\mathcal{L}_{\mathrm{MSE}}+10^{-4}\mathcal{L}_{\mathrm{KL}}$};

\draw[arr] (x)--(enc); \draw[arr] (enc)--(mu); \draw[arr] (enc)--(lv);
\draw[arr] (mu)--(z); \draw[arr] (lv)--(z);
\draw[arr] (z)--(dec); \draw[arr] (z)--(head);
\draw[arr] (dec)--(rec); \draw[arr] (head)--(prob);
\draw[arr] (rec)--(loss); \draw[arr] (prob)--(loss);
\end{tikzpicture}

\caption{Variational multi-task architecture used as the basis of the breed-aware genomic representation model. The encoder parameterizes a latent posterior regularized by a KL-divergence term; the resulting latent variable supports both genotype reconstruction and breed prediction. Author's adaptation of the VAE formulation \cite{kingma2013auto} and the VMGP multi-task paradigm \cite{zhao2025vmgp}.}
\label{fig:vmgp_architecture}
\end{figure}



\section{Hyperspherical Contrastive Representation Model}
\label{sec:contrastive_model}

The second representation-learning approach follows a contrastive
learning objective \cite{oord2018representation,chen2020simple} and
is considered as a complementary alternative to the variational model.

The motivation for this choice is that accurate reconstruction of the
complete genotype profile is not necessarily required to obtain a
useful representation for classification. A generative model may
preserve genomic variation that is not directly relevant to the
downstream task. Contrastive learning introduces a different
inductive bias: rather than reconstructing every input variable, it
learns the representation space through similarity relationships
between different views of the observations.

In the present setting, stochastic transformations of the same
genotype profile are treated as related views. The model is therefore
encouraged to produce similar representations for these views while
separating them from other observations. This objective is intended
to promote genomic features that remain stable under the adopted
perturbations, without requiring explicit reconstruction of the full
SNP vector.

The contrastive representations are additionally normalized to the
unit hypersphere, where the dot product between two representations
corresponds to cosine similarity
\cite{wang2020understanding}. This geometry provides a natural space
for similarity-based optimization.

The variational and contrastive models therefore represent two
complementary hypotheses about genomic representation learning:
the former emphasizes regularized compression together with
reconstruction and breed prediction, whereas the latter organizes
the latent space according to similarity between related views.
Their comparison under the same experimental protocol allows these
different inductive biases to be evaluated empirically.

For each genotype vector $\mathbf{x}_i$, a stochastic augmentation
operator produces two correlated views

\begin{equation}
    \tilde{\mathbf{x}}_i^{a},
    \tilde{\mathbf{x}}_i^{b}
    \sim
    \mathcal{T}(\mathbf{x}_i),
    \label{eq:contrastive_views}
\end{equation}

where $\mathcal{T}$ denotes the genomic augmentation procedure.

% TODO:
% Describe exactly what T does.
% If the augmentation is designed specifically for this thesis, say so.
% If it is adopted from a genomic paper, cite that exact source.

Both views are processed by a shared encoder $f_{\theta}$ and
projection head $g_{\psi}$:

\begin{equation}
    \mathbf{h}_i^{a}
    =
    g_{\psi}
    \left(
        f_{\theta}
        (
            \tilde{\mathbf{x}}_i^{a}
        )
    \right),
\end{equation}

\begin{equation}
    \mathbf{h}_i^{b}
    =
    g_{\psi}
    \left(
        f_{\theta}
        (
            \tilde{\mathbf{x}}_i^{b}
        )
    \right).
\end{equation}

The projection vectors are normalized as

\begin{equation}
    \mathbf{z}_i
    =
    \frac{
        \mathbf{h}_i
    }{
        \|
            \mathbf{h}_i
        \|_2
    },
    \qquad
    \|
        \mathbf{z}_i
    \|_2
    =
    1.
    \label{eq:hyperspherical_normalization}
\end{equation}

Normalized contrastive representations can be interpreted
geometrically on a hypersphere, where the dot product of unit
vectors corresponds to cosine similarity
\cite{wang2020understanding}.

Training uses the NT-Xent objective introduced in the SimCLR
framework
\cite{chen2020simple}:

\begin{equation}
\begin{split}
    \mathcal{L}_{\mathrm{NT-Xent}}(i,j)
    =
    -\log
    \frac{
        \exp
        \left(
            \operatorname{sim}
            (
                \mathbf{z}_i,
                \mathbf{z}_j
            )
            /
            \tau
        \right)
    }{
        \displaystyle
        \sum_{k=1}^{2N}
        \mathbb{I}_{[k \neq i]}
        \exp
        \left(
            \operatorname{sim}
            (
                \mathbf{z}_i,
                \mathbf{z}_k
            )
            /
            \tau
        \right)
    },
\end{split}
\label{eq:nt_xent}
\end{equation}

where $\tau$ denotes the temperature parameter.

% TODO (optional figure): add the final contrastive architecture diagram
% after the augmentation operator and implemented layers are finalized.




The implemented contrastive architecture is illustrated in
Figure~\ref{fig:contrastive_architecture}.

\begin{figure}[htbp]
\centering
\begin{tikzpicture}[
    >=Latex,
    font=\small,
    box/.style={draw, semithick, rounded corners=2pt, align=center,
                minimum height=10mm, minimum width=37mm, inner sep=4pt, fill=black!2},
    key/.style={box, fill=black!7},
    small/.style={draw, semithick, rounded corners=2pt, align=center,
                  minimum height=9mm, minimum width=31mm, inner sep=3pt, fill=black!2},
    arr/.style={->, semithick}
]
\node[key] (x) at (0,4.65) {Genotype profile $\mathbf{x}_i$};
\node[box] (aa) at (-3.25,3.15) {Augmented view A};
\node[box] (ab) at (3.25,3.15) {Augmented view B};
\node[box] (ea) at (-3.25,1.65) {Shared encoder $f_\theta$};
\node[box] (eb) at (3.25,1.65) {Shared encoder $f_\theta$};
\node[small] (za) at (-3.25,0.15) {$L_2$ norm. $\rightarrow \mathbf{z}_i^{a}$};
\node[small] (zb) at (3.25,0.15) {$L_2$ norm. $\rightarrow \mathbf{z}_i^{b}$};
\node[key] (np) at (0,-1.35) {Centroid N-pair loss (Eqs.~3--6)};
\node[key] (sphere) at (0,-2.90) {Hyperspherical latent space $\mathbb{S}^{2}$};

\draw[arr] (x) -- (aa); \draw[arr] (x) -- (ab);
\draw[arr] (aa) -- (ea); \draw[arr] (ab) -- (eb);
\draw[arr] (ea) -- (za); \draw[arr] (eb) -- (zb);
\draw[<->, semithick] (za) -- node[above]{positive pair} (zb);
\draw[arr] (za) -- (np); \draw[arr] (zb) -- (np);
\draw[arr] (np) -- (sphere);
\end{tikzpicture}

\caption{Self-supervised hyperspherical pipeline. Two independently
augmented views of the same individual are encoded by a shared
convolutional encoder and $L_2$-normalized onto the unit sphere. The
representation is optimized with the centroid N-pair objective, without
a projection head, without NT-Xent temperature and without breed labels
\cite{thor2025contrastive,wang2020understanding}.}
\label{fig:contrastive_architecture}
\end{figure}



\section{Sample-Size Evaluation Protocol}
\label{sec:sample_size_protocol}

To address \textbf{RQ2}, the number of training individuals available
per breed is varied systematically.

The decision to treat sample size as an experimental factor is also
motivated by previous livestock-genomics work showing that the
relative behavior of genomic analysis methods can change
substantially as the number of available animals increases
\cite{manca2020multivariate}.
That study concerns genotype--phenotype association rather than
breed classification and therefore does not provide a numerical
sample-size threshold for the present task; it is used here only as
methodological motivation.

Let

\begin{equation}
    \mathcal{N}
    =
    \{
        n_1,
        n_2,
        \ldots,
        n_K
    \}
    \label{eq:sample_sizes}
\end{equation}

denote the set of evaluated per-breed training sample sizes.

For every $n_k$, individuals are sampled from the training partition
while preserving the held-out validation and test partitions.

Repeated sampling with independent random seeds is used to reduce the
dependence of the analysis on a particular subset of individuals.

% TODO:
% Insert the actual values only once decided.
%
% Example:
% \mathcal{N} = \{5, 10, 15, 20, 30, 40, 50\}
%
% Do NOT write that 30 is the minimum here.

For each sample-size configuration, classification performance is
summarized across repeated runs using both a central estimate and a
measure of variability.

% TODO: add the exact Results-section cross-reference once the
% sample-efficiency subsection and its label are finalized.


\section{Evaluation Metrics}
\label{sec:evaluation_metrics}

Breed-classification performance is evaluated using multiple
complementary metrics.

Standard accuracy is reported together with macro-averaged F1 score
and balanced accuracy.
Macro-averaged measures are useful in multiclass evaluation because
they aggregate class-specific performance while assigning equal
weight to each class
\cite{sokolova2009systematic}.

For $C$ breeds,

\begin{equation}
    F_{1}^{\mathrm{macro}}
    =
    \frac{1}{C}
    \sum_{c=1}^{C}
    F_{1,c}.
    \label{eq:macro_f1}
\end{equation}

Confusion matrices are additionally used to inspect which breed
pairs are systematically difficult to distinguish.

Where experiments are repeated across multiple random seeds,
performance is summarized using a central estimate together with an
explicit measure of variability.

% TODO: choose and report the final uncertainty measure used throughout
% the experiments (e.g., standard deviation, standard error, or a
% confidence interval). Once chosen, define its notation here.


\section{Operational Definition of Minimum Requirements}
\label{sec:minimum_requirements}

The expressions ``minimum number of samples'' and ``minimum number
of markers'' require an explicit operational definition.

For marker reduction, let

\begin{equation}
    S_{\mathrm{full}}
\end{equation}

denote the classification score obtained using the complete
post-QC SNP panel.

For a panel containing $P'$ selected markers, let

\begin{equation}
    S(P')
\end{equation}

denote the corresponding downstream classification score.

The minimum retained panel is operationally defined as

\begin{equation}
    P_{\min}
    =
    \min
    \left\{
        P'
        :
        S(P')
        \geq
        (1-\epsilon_P)
        S_{\mathrm{full}}
    \right\}.
    \label{eq:minimum_marker_definition}
\end{equation}

Similarly, if $S(N')$ denotes performance obtained using $N'$
training individuals per breed, the minimum sample requirement is analogously defined as

\begin{equation}
    N_{\min}
    =
    \min
    \left\{
        N'
        :
        S(N')
        \geq
        (1-\epsilon_N)
        S_{\mathrm{reference}}
    \right\}.
    \label{eq:minimum_sample_definition}
\end{equation}

The tolerances $\epsilon_P$ and $\epsilon_N$ must be fixed before
drawing the final conclusions from the corresponding curves.

% TODO: fix the operational tolerances before inspecting the final
% conclusions. Record the agreed criterion here, e.g.:
% - 95% of full/reference performance;
% - 99% of full/reference performance;
% - within 1 percentage point of the reference;
% - statistically indistinguishable from plateau performance.

The concrete panel-construction procedure is described in
Chapter~\ref{chap:feature_importance}, while the resulting
sample-size and marker-size curves are reported in
Chapter~\ref{chap:experiments}.
